\subsection{Data Collection}

Our data come from a single Unitree Go2, whose onboard proprioception is published on the ROS~2 topic \texttt{/lowstate}~\cite{macenski2022ros2,unitree_go2}. We recorded this stream while running scripted locomotion scenarios and stored each as a comma-separated-value (CSV) log. Positive scenarios walk the robot and then drive one leg into a physical entanglement (a wire or a hand) before releasing it; negative scenarios contain no entanglement. We also scripted a \texttt{Lock} condition, the rigid stance the robot holds right after standing up, because it produces large, sustained joint torques with almost no motion and thus superficially resembles a snag; including it as a labeled negative teaches the model to tell a held stance from a genuine entanglement.

Logging is event-driven, so the effective rate is non-uniform, ranging from roughly $410$ to $500$~Hz (median near $483$~Hz); we keep this raw timing and defer resampling to the feature stage. The corpus is $25$ recordings totaling $159{,}001$ labeled rows, each with $51$ columns in three sensor groups plus a label. The motor group holds $36$ values: four legs, each with hip, thigh, and calf joints reporting position $q$, velocity $dq$, and torque $\tau$. The next group is the four foot-contact forces, and the last is the inertial measurement unit (roll, pitch, yaw, three gyroscope and three accelerometer axes). A final \texttt{Status} column carries the phase label. Table~\ref{tab:dataset} summarizes these groups.

We labeled the recordings by hand. Because each scenario was scripted, we knew when an entanglement began and ended, and we combined those timestamps with per-joint torque plots (the more reliable cue, since a snag shows a distinctive sustained resistive torque) to mark the exact rows of each event. Each row receives one of the phase labels \texttt{Walking}, \texttt{Entangled}, \texttt{Stop}, \texttt{Lock}, or blank, and the affected leg of a positive recording is encoded in its filename.

Of the $25$ recordings, $15$ are positive (one contiguous entangled event each) and $10$ are negative. The positive events are unevenly distributed across legs: grouping by affected leg gives seven recordings for the rear-right, five for the rear-left, four for the front-left, and only three for the front-right. The front-right leg is therefore under-represented, which constrains how confidently its per-leg behavior can be estimated and resurfaces in the per-leg results (Table~\ref{tab:perleg}). Splits are formed at the level of whole recordings (Sec.~\ref{sec:method} explains why), assigning $16$ recordings to training, $4$ to validation, and $5$ to test; the held-out test set is \texttt{back\_left\_hand2}, \texttt{back\_right\_wire1}, \texttt{front\_left\_hand2}, and \texttt{front\_both\_wire2}, with \texttt{go2\_lowstate\_lock\_stop} as the negative.
